\section{Conclusion and extensions}

\subsection{Conclusion}

The research question asked how the separation of the magnets affects the two normal-mode frequencies of a coupled magnetic-mechanical oscillator, and hence its average and beat frequencies; the quantitative test adopted was whether the coupling term $f_2^2-f_1^2=2f_{\rm beat}f_{\rm avg}$ follows the inverse-fifth-power law of the point-dipole model.

The answer to the first part is clear and agrees with the theory of coupled oscillators. The in-phase frequency $f_1\approx0.81\un{Hz}$ is almost independent of the separation (it varies by $2.6\,\%$ over $21$--$40\un{mm}$, a variation that Section~\ref{sec:evaluation} attributes partly to the analysis and partly to a small difference between the springs), while the anti-phase frequency $f_2$ falls monotonically from $1.063\un{Hz}$ at $21\un{mm}$ towards $f_1$, the two becoming indistinguishable in a $30\un{s}$ record beyond $40\un{mm}$. The magnetic repulsion behaves as a contactless coupling spring whose stiffness, relative to the leaf springs, can be tuned from about $0.7$ to below $0.02$ by moving the clamps a few centimetres. Consequently the average frequency stays within $1\,\%$ of $f_1$ beyond $30\un{mm}$ and rises by $16\,\%$ at the closest separation, while the beat frequency, which is the mode splitting itself, falls from $\nfbeatSmall\un{Hz}$ at $21\un{mm}$ to below the resolution of a $30\un{s}$ record ($33\un{mHz}$) beyond $40\un{mm}$: the beats slow from one every $\nTbeatSmall\un{s}$ to one every few minutes.

The answer to the second part is more nuanced. Over $21$--$35\un{mm}$ the coupling term follows a power law well ($R^2=0.989$), but with exponent $-3.75\pm0.19$ (unweighted) or $-3.5\pm0.2$ (weighted or non-linear fit) rather than $-5$. The evaluation traced this $25$--$30\,\%$ shortfall to a combination of effects, each with a definite sign. Reading two nearby peaks from the spectrum of a $30\un{s}$ record biases $f_2-f_1$ upwards at large $d$ and accounts for roughly one third of the shortfall, as a simulation of the complete measurement chain shows. The finite size of the magnets, which are only $2$--$6$ diameters apart, makes the true magnetic stiffness fall more slowly than $d^{-5}$ in this range, and a model with uniformly magnetised cylinders instead of point dipoles fits the data as well as a free power law, with a diameter of $\nDfit\pm\nDfitErr\un{mm}$. A small difference between the two springs may contribute at the large-$d$ end and explains the drift of $f_1$. The one effect that acts in the opposite direction, the hardening of the magnetic force at the $10\un{mm}$ release amplitude, steepens the line by about $0.15$ and therefore slightly masks the others. The point-dipole law is not contradicted by these results; it is the large-separation limit of a force law whose short-range form is measurably different, and the data are consistent with that picture.

The main methodological result of the essay is that the limiting factor was not the apparatus but the analysis. Tracking both magnets and Fourier-transforming the normal coordinates $x_1\pm x_2$ removes the two-peak resolution problem entirely, and a simulation shows that it recovers the exponent to $\pm0.01$ over $21$--$60\un{mm}$ where the original method fails beyond $35\un{mm}$. Together with a smaller release amplitude, longer records and a finite-size force law, this gives a measurement capable of distinguishing the remaining explanations. \TBM{When the re-measurement has been done: state the final exponent, its uncertainty, and the range of $d$ over which the dipole law holds to within the uncertainty.}

\subsection{Extensions}

\begin{itemize}
\item \textbf{Test the $\mum^2$ dependence.} Repeat at fixed $d$ with magnets of different grade (N35, N42, N52) or stacked pairs: equation~\eqref{eq:main} predicts $\Y\propto\mum^2$, a test of the magnitude of the coupling that is independent of the distance law.
\item \textbf{Attractive orientation.} With unlike poles facing, the magnetic stiffness changes sign and \emph{softens} the anti-phase mode: $\omega_2^2=(k-2|\gamma|)/m$. As $d$ is reduced, $f_2$ should fall to zero at a critical separation where the springs snap together, a buckling-type instability whose position tests the same force law from the other side.
\item \textbf{Three or more springs.} A chain of $N$ magnet-tipped springs has $N$ normal modes with frequencies given by a tridiagonal matrix; with the tunable magnetic coupling it is a table-top model of a one-dimensional lattice, and the dispersion relation could be mapped.
\item \textbf{Eddy-current damping.} Placing a copper or aluminium plate behind the magnets adds a velocity-dependent damping that can be varied with the plate distance, allowing the transition from beats to over-damped energy transfer to be studied.
\item \textbf{The equilibrium quintic.} Measuring the static outward deflection $\delta$ against the clamp separation tests the same $r^{-4}$ law statically and would give $\mum$ from the same apparatus without a balance.
\end{itemize}
